% ЕМТ 2022, VIII клас — точен LaTeX препис от официалните файлове в Klasirane.com.
% Условия: https://klasirane.com/api/competitions/EMT/All/2022/8%20%D0%BA%D0%BB/probs
% SHA-256 (условия): 176dd0576679c73d67aa1104744f1969197e5325256235b83189ad1dd52bbf8a
% Решения: https://klasirane.com/api/competitions/EMT/All/2022/8%20%D0%BA%D0%BB/sol
% SHA-256 (решения): 4a2efd0afb1c16514a04f58f3019bd6b3b5c01d45c096413add97567bb58e15c
% Mathpix Snip (условия): 53817608-aa05-4ab7-b2db-9a86949650d4
% Mathpix Snip (решения): 24797e8d-c3c6-46ac-9c50-83a1b9cfee8c
% Точковите оценки, правописът, формулировките и видимите печатни особености са запазени от източника.
% В официалните файлове няма фигури към задачите за VIII клас.

Задача 8.1. Решете уравнението
$$
4 x^{2}+|9-6 x|=|10 x-15|+6(2 x+1) .
$$

Отговор. $x=4, x=-1$.

Решение. Първи метод. Преобразуваме еквивалентно:
$$
\begin{gathered}
4 x^{2}-6(2 x+1)=5|2 x-3|-3|3-2 x| \\
2 x^{2}-6 x-3=|2 x-3|
\end{gathered}
$$
Понеже $|a|$ е винаги $a$ или $(-a)$, то решенията изпълняват $2 x^{2}-6 x-3=2 x-3$ или $2 x^{2}-6 x-3=3-2 x$. В първия случай следва $2 x(x-4)=0$, т.е. $x=0$ или $x=4$, а във втория случай $2(x-3)(x+1)=0$, т.е. $x=3$ или $x=-1$. Обаче $x=0$ и $x=3$ дават в началното уравнение $-3=3$, което е невярно, докато $x=-1$ и $x=4$ наистина са решения (и двете страни са равни на 5).

Втори метод. Преобразуваме еквивалентно до
$$
4 x^{2}-12 x-6=2|2 x-3| .
$$
Полагаме $u=|2 x-3| \geq 0$; тогава $u^{2}=4 x^{2}-12 x+9$ и уравнението добива вида $u^{2}-15=2 u$, което води до $(u-5)(u+3)=0$. Вторият множител е положителен, така че трябва $u=5$. Съответно $2 x-3=5$ води до $x=4$, а $2 x-3=-5$ води до $x=-1$.

Оценяване. (6 точки) 4 т. за намиране на 4, $(-1)$ и още най-много две други $x$, 2 т. за отхвърляне на излишните стойности.

Задача 8.2. Даден е правоъгълен триъгълник $A B C$ с прав ъгъл при върха $C$ и лице $S$. Нека $S_{1}$ е лицето на кръга с диаметър $A B$ и $k=\frac{S_{1}}{S}$.

а) Да се намерят острите ъгли на $A B C$, ако $k=2 \pi$.

б) Да се докаже, че не съществува $A B C$, за който $k=3$.

Решение. Ако $A B=c$ и $h$ е височината към $A B$, то $S=\frac{c h}{2}, S_{1}=\frac{\pi c^{2}}{4}$ и $k=\frac{\pi}{2} \cdot \frac{c}{h}$. Нека $m=\frac{c}{2}$ е дължината на медианата към хипотенузата.

а) При $k=2 \pi$ следват $c=4 h$ и $m=2 h$. От правоъгълния триъгълник с хипотенуза $m$ и катет $h$ следва, че ъгълът между медианата и хипотенузата е 30°, следователно острите ъгли на $A B C$ са $15^{\circ}$ и $75^{\circ}$.

б) От правоъгълния триъгълник с хипотенуза $m$ и катет $h$ имаме $m \geq h$. Така $\frac{c}{h} \geq 2$ и значи $k \geq \frac{\pi}{2} \cdot 2=\pi>3$.

Оценяване. (6 точки) 2 т. за а), от които 1 т. за $c=4 h$ и 1 т. за довършване (не се отнемат точки, ако свойството на правоъгълен триъгълник с ъгъл 15° се цитира като известен факт); 4 т. за б), от които: 1 т. за изразяване на $k$ чрез $c$ и $h$, 1 т. за въвеждането на $m$, 1 т. за $m \geq h$ и 1 т. за довършване; алтернативно: 1 т. за изразяване на $k$ чрез катетите $a$ и $b$ на $A B C$, 1 т. за твърдението $\frac{a^{2}+b^{2}}{a b} \geq 2$ и 1 т. за доказателство, 1 т. за довършване.

Задача 8.3. По окръжност са разположени в този ред точките $A_{1}, B_{1}, A_{2}, B_{2}, \ldots, A_{9}, B_{9}$. Всяка от отсечките $A_{i} B_{j}(i, j=1,2, \ldots, 9)$ трябва да се оцвети в един от $k$ дадени цвята, така че никои две едноцветни отсечки не се пресичат във вътрешна точка и за всяко $i=1, \ldots, 9$ има цвят, за който в този цвят няма отсечки с край $A_{i}$, нито с край $B_{i}$. Намерете най-малкото възможно $k$.

Отговор. 9

Решение. Отсечките $A_{1} B_{5}, A_{2} B_{6}, A_{3} B_{7}, A_{4} B_{8}, A_{5} B_{9}, A_{6} B_{1}, A_{7} B_{2}, A_{8} B_{3}$ и $A_{9} B_{4}$ се пресичат във вътрешни точки, така че са необходими поне 9 цвята. Толкова са и достатъчни: може в цвят $i$ да са отсечката $A_{i+1} B_{i}$ (навсякъде при $i=9$ пишем 1 вместо $i+1$) и всички отсечки с край $A_{i}$ освен $A_{i} B_{i-1}$ (при $i=1$ пишем 9 вместо $i-1$). По този начин никоя от отсечките с край $A_{i}$ и/или $B_{i}$ не е в цвят $i+1$.

Забележка. Дадената задача е формализация на следната постановка: на кръстовище излизат 9 двупосочни улици с дясно движение; всяка улица се пресича от пешеходна пътека точно до кръстовището. Има светофарна уредба, осигуряваща зелен сигнал за преминаване от всяка улица по най-краткия път към всяка улица, включително обратен завой към самата нея, както и за преминаване по всяка пешеходна пътека. Такт наричаме период от време, за който множество от зелени сигнали за автомобили и/или пешеходци не се променя и маршрутите им не се пресичат. Намерете най-малкия възможен брой тактове (всяко платно може да бъде разделено на достатъчен брой ленти за автомобилите отиващи/идващи към/от различни улици).

Забележка. По подобен начин (с лека модификация на доказателството на оценката при четно $n$) можем да докажем, че отговорът в горната задача в случая на $n$ улици ($n \geq 3$) е $n$.

Оценяване. (7 точки) 3 т. за доказателство, че са необходими поне 9 цвята, 4 т. за работеща схема с 9 цвята.

Задача 8.4. Да се намери броя на редиците от 2022 естествени числа, такива че във всяка редица:

\begin{itemize}
\item всяко число след първото е по-голямо или равно на предходното,
\item поне едно от числата е равно на 2022 и
\item сумата на всеки 2020 от числата се дели на всяко от останалите две.
\end{itemize}

Отговор. 13

Решение. Нека сумата на числата е $S$ и $a$, $b$ и $c$ са кои да е три от тях. Явно $S-a-b$ и $S-a-c$ се делят на $a$, откъдето $a$ дели $b-c$. Сега ако изберем $a$ да е число с най-голяма стойност измежду всички, неравенството $a \leq|b-c|$ няма как да бъде изпълнено и така горната делимост дава непременно $b=c$. Следователно всяка от търсените 2022-орки непременно има вида $(n, n, \ldots, n, a)$. Вече исканото е еквивалентно на делимостите $a \mid 2020 n$ и $n \mid 2019 n+a$, т.е. $n \mid a \mid 2020 n$ – значи търсените са от вида $(n, n, \ldots, n, k n)$, където $k$ е делител на 2020 и поне едно от $n$ и $k n$ е равно на 2022.

Ако $n=2022$, то понеже $2020=2^{2} \cdot 5 \cdot 101$ има $(2+1) \cdot(1+1) \cdot(1+1)=12$ делителя, имаме 12 възможности.

Ако $k n=2022$ и $n \neq 2022$, то $k>1$ дели 2020 и 2022, значи дели 2, т.е. единствената възможност е $k=2$ и $n=1011$.

Оценяване. (7 точки) 1 т. за верен отговор; 2 т. за $a \mid b-c$ за произволни $a$, $b$, $c$; 1 т. за разглеждане на максимален елемент; 1 т. за свеждането до вида $(n, n, \ldots, n, k n)$ за $k \mid 2020$; по 1 т. за $n=2022$ и $k n=2022$.
